\subsection{I.3 Le projet d'une indexation automatisée des archives audiovisuelles, conditionnée par le modèle en place}

le catalogage et l'indexation des archives audiovisuelles est une opération laborieuse et longue, particulièremet sur des archives spécifiques telles que nous vons décrites 

\subsubsection{Projet SKAI}
La mise en place de l'IA est justifiée par des éléments que nous avons déjà évoqué : les marqueurs temporels. Elle répond à un besoin concret d'un corpus précis, mais pourra être étendu à toutes les institutions préservatrices de collections audiovisuelles. 

réponse aux problèmes du modèle en place ? pas vraiment, ajout e nouvelles fonctionnalités pour expérience utilisateur globale

parmi les fonctionnalités évoquées : indexation audiovisuelle > segmentation temporelle 

issue d'une demande client, marqueurs temporels déjà présents dans les données. La segmentation viendrait automatiser cette pratique et l'étendre à des archives non annotées

répondre à des besins utilisateurs concrets et non pas seulement à une logique produit 


- moteurs IA 
- objectif de transparence 
- clients qui refusent l'IA 
- projets en cours 
- répondre à un besoin utilisateur et non seulement à une logique produit (tension public/privé) (voir fichier contraintes institution)

"Alors que l'utilisation du ML	va déjà de soi dans certains domaines, il y a des domaines dans lesquels nous devons	rendre son utilisation transparente pour les utilisateur·rice·s et leur donner la possibilité
d'option de retrait, par exemple lorsqu'il s'agit de données personnelles. La transparence est également requise lorsqu'il s'agit d'adapter l'offre concrète aux besoins des utilisateur·rice·s	sur la base de profils d'utilisateur·rice·s (obtenus par le biais de demandes, d'autodéclaration ou de « traces d'utilisation », comme le classement d'une liste de résultats ou	des propositions par des systèmes de recommandation)."
\footcite{zotero-item-584}

ce projet doit être intégré dans une pratique d'indexation, nécessaire à l'intéropérbilité des données et au bon foctionnement du modèle. 

\subsubsection{Indexation}
définition > finalité pour l'expérience utilisateur > a pour condition la segmentation et la description 

définition de l'indexation :
"L'indexation est cette technique consistant à caractériser le contenu d'un document et l'information qu'il détient de manière à le retrouver quand on effectue des recherches sur l'un des sujets dont il traite. La difficulté est donc de savoir caractériser et représenter l'information documentaire pour qu'il soit aisé de la mettre en rapport avec des sujets d'investigation. Mise en rapport d'une requête et d'un contenu représenté et synthétisé, l'indexation permet de s'orienter dans la masse des documents et d'organiser ses connaissances. L'indexation appartient donc à ce qu'on appelle depuis quelques années les techniques intellectuelles Indexation et archivage de contenus multimédias"\footcite{zotero-item-582}

Définition : "Elle est tout d abord la condition préalable pour fournir un point d accès sûr et vérifié à l information en utilisant des notions clairement définies, représentées par des termes eux-mêmes soigneusement choisis et contrôlés, formant des vocabulaires contrôlés (listes d autorité, thésaurus) ou des notices d autorité publics, accessibles à tout un chacun et maintenus, complémentaires du langage « naturel » des descriptions archivistiques, voire des documents eux-mêmes." \footcite{zotero-item-675}

"En permettant d affiner les recherches, elle améliore la
découvrabilité des données. Elle rend possible la création de formulaires pertinents, base pour la construction d accès plus sophistiqués : modes de navigation intuitifs dans les instruments de recherche (filtres, facettes), nouveaux services reposant sur
la visualisation de données (à commencer par l accès géographique tant réclamé par le public). (cf améliorations s movies + recherche)"\footcite{zotero-item-675}

"l'indexation est un outil de partage et de publication des données sur le web" \footcite{zotero-item-675}

L'indexation se fait en deux étapes : segmentatin et description du contenu 

l'indexation manuelle est u travail long et laborieux étant donn la nature de l'image en mouvement 

\subsubsection{Segmentation}
segmentation virtuelle vs physique > les différentes segmentations possibles en audiovisuel > 

La segmentation appliquée à l'image en général, y compris l'image fixe, consiste à diviser l'image en groupes de pixels selon certains critères. 

définition de la segmentation auiovisuelle : 
"Il s’agit d’un découpage conceptuel et non d’un montage technique: le fichier vidéo n’est pas modifié physiquement ; seules des balises temporelles délimitent les segments retenus.
Cette approche garantit la souplesse de l’analyse, permet des lectures multiples d’un même document et rend possible la production de descriptions structurées" \footcite{stockinger2011}

"la segmentation ne se limite pas à un simple acte technique : elle représente un moment clé de la transformation d’un document audiovisuel" \footcite{stockinger2011}

"La segmentation est définie comme une étape préalable à toute analyse, qui consiste à délimiter des sous-ensembles homogènes à l’intérieur d’un texte audiovisuel." \footcite{stockinger2011}

la segmentation est étroitement liée à l'interprétation

la segmentation est issue d'une problématique qui pose le contexte et l'objectif donné. Elle définit également le type d'action que permet la segmentation : par exemple l'indexation, la description, l'analyse... \footcite{stockinger2011}

Il faut aussi définir la nature de la segmentation, sa granularité, la nature des segments, le niveau de description attendu \footcite{stockinger2011}

Une distinction est également établie entre "segmentation virtuelle", le fait de poser des repères sur l'image pour en déceler les séquences pertinentes, et le "découpage physique" qui crée de nouveaux documents, scindés depuis l'archive originale. Cela permettrait de mettre en logne des passages autonomes d'une oeuvre par exemple. \footcite{stockinger2011}

"La segmentation agit comme un filtre de pertinence : elle configure le champ d’observation en sélectionnant ce qui deviendra objet d’analyse sémiotique, d’indexation ou de publication. [...] segmenter, c’est déjà choisir un point de vue sur le document." \footcite{stockinger2011}


"la segmentation apparaît comme un premier niveau de
modélisation, à la frontière entre observation empirique et structuration conceptuelle. Elle engage une série de décisions qui feront de la vidéo non plus seulement un enregistrement, mais un objet de connaissance" \footcite{stockinger2011}

les perspectives de segmentation souhaitées : "la possibilité d’exercer plusieurs segmentations parallèles (visuelle,
sonore, linguistique) sur un même fichier" \footcite{stockinger2011}

Appliquée à la vidéo, la segmentation peut alors être de deux formes : une division en plusieurs segments ou objets. \footcite{zhou2022}. La division d'une image animée en séquences ou en différents plans correspond à une segmentation temporelle. La segmentation d'objets dominants dans l'image \footcite{zhou2022}
selon des critères visuels\footcite{carrive2014}, désigne plutôt une segmentation spatiale, analysant la place des objets dans l'image, pas encore leur signification.\footcite{zhou2022}. 
C'est au moment d'apposer une signification sur ces objets qu'on parle de segmentation sémantique. En effet, la segmentation sémantique consiste à extraire des objets de l'image selon des catégories sémantiques prédefinies. \footcite{zhou2022} Ce type de segmentation prépare l'indexation en assignant des étiquettes à chaque pixel extrait de l'image. 
Parmi ces trois types de segmentation, seule la segmentation temporelle est spécifique à l'image animée, puisqu'elle analyse la différence temporelle et sémantique entre les plans et leur relations pour pouvoir en extraire des séquences distnctes. Les segmentation d'objets et sémantiques sont en revanche également applicables à l'image fixe. 

Enfin, peut être appliquée à une vidéo une segmentation de la bande sonore, qu'on désignera sous le nom de segmentation audio, utilisée par exemple, pour segmenter les paroles de la musique das une même séquence vidéo. \footcite{carrive2014}